\chapter{实验设置与结果分析}

\section{实验环境与配置}

实验在 PyTorch 框架下实现，训练脚本按阶段独立执行并统一记录结构化日志。
核心设计参数包括焦距、F 数、镜片厚度预算与衍射面参数化方式。

\section{阶段结果与对比}

折射阶段输出结构参数与 MTF/Spot 结果，折衍阶段输出相位分布与损失收敛曲线，端到端阶段输出重建质量指标。
多次重复实验用于验证结果稳定性，并在统一随机种子管理下报告均值与方差。

\section{误差来源与边界讨论}

当前误差主要来自离轴场点传播近似、参数初始化敏感性与数据分布偏移。
后续将进一步引入制造约束和多场景验证，以提升结论在真实系统中的泛化能力。
